\section{Methodology}

\noindent Excel-Shrink is a stream-oriented XLSX cleaner that removes visually inconsequential XML
while preserving user-visible layout semantics
(e.g., fill/border styles that convey meaning),
and merges superfluous column definitions.
The system operates directly on the bytestream of the zipped OOXML container:
it parses \texttt{xl/workbook.xml}
and per-sheet \texttt{xl/worksheets/sheet*.xml} streams with regular expressions
instead of loading full XML trees into memory,
applies workbook-specific rewriting rules
derived from \texttt{xl/styles.xml} and \texttt{xl/sharedStrings.xml},
and re-emits a compacted archive.

We first formalise inputs and artefacts, then describe pipeline stages.

\subsection{Inputs, Outputs, and Notation}

\textbf{Inputs.}
\emph{Path} may be a single XLSX file, a directory, or an HTTPS/HTTP URL. Optional flags control recursion, multiprocessing, chunk size for streaming, sheet filters, and whether only the workbook is cleaned.

\textbf{Outputs.}
A cleaned XLSX written either to a file path or into an output directory, preserving all unmodified parts of the original OOXML.

\textbf{Core artefacts.}
\begin{itemize}
\item \emph{Workbook XML} (\texttt{xl/workbook.xml}): contains \texttt{<definedNames>} nodes.
\item \emph{Sheet XML} (\texttt{xl/worksheets/sheet*.xml}): contains \texttt{<cols>}, \texttt{<sheetData>}, \texttt{<mergeCells>}.
\item \emph{Styles} (\texttt{xl/styles.xml}): provides \texttt{<cellXfs>} with fill/border/apply flags.
\item \emph{Shared Strings} (\texttt{xl/sharedStrings.xml}): enables identification of pure-whitespace string table entries.
\end{itemize}

\noindent \textbf{Workbook-specific patterns.}
From \texttt{styles.xml} and \texttt{sharedStrings.xml} we derive:
\begin{enumerate}
\item sets of style indices with \emph{fill} or \emph{border} or with neither but \emph{font-only};
\item regex-like predicates for: (i) whitespace cells, (ii) consecutive empty cells/rows, (iii) attributes irrelevant for empty rows (\emph{spans}, vendor \emph{dyDescent}, etc.), and (iv) presence of \emph{hidden} rows/columns.
\end{enumerate}





\subsection{Deriving Workbook-Specific Patterns}
\label{sec:wb-patterns}


Our worksheet cleaning approach adopts a two-stage cleaning procedure based on a chunking strategy.
In the initial pre-processing phase, individual chunks are analysed
and cleaned using lightweight operations that do not require any global context
from the worksheet. Some of these operations require knowledge of the workbook structure,
others can be applied independently to each worksheet. Despite their simplicity,
these early operations are effective in eliminating a large portion of file bloat
from Excel spreadsheets.

Such operations involve the removal of visually inconsequential cells and rows—
cells that contain no content and no layout-relevant styling,
such as background fills or borders.
Font-only styling (e.g., font family or size) or filled  with
 whitespace characters  is insufficient to classify a cell as visually relevant.
By filtering out these cells during chunk-level processing,
we reduce both the computational burden and the memory footprint
of subsequent operations.

To identify and remove these visually inconsequential cells,
we need to know which styles correspond to visually empty formatting
and which shared strings correspond to whitespace-only content.

After extracting the style definitions from styles.xml
and the shared strings from sharedStrings.xml during the initial context extraction phase,
we can compile regex patterns that match cells and rows that are visually empty.
The  compilation of these patterns is summarised in Algorithm \ref{alg:build-wb-patterns}.

\begin{algorithm}
\DontPrintSemicolon
\SetKwFunction{ParseCellXfs}{parseCellXfs}
\SetKwFunction{IdxWhitespaceSI}{indicesOfWhitespaceSI}
\SetKwFunction{Join}{join} % textual join with "|"
\SetKwProg{Fn}{Function}{:}{}

\Fn{BuildWorkbookPatterns(stylesXML, sharedStringsXML)}{

  \tcp{1) Parse XF records}
  $xfs \leftarrow \ParseCellXfs(\text{stylesXML})$\;

  \tcp{2) Partition indices by properties}
  $S_{\mathrm{fb}} \leftarrow \{\, i \mid xfs[i].\mathrm{fillId} \neq 0 \lor xfs[i].\mathrm{borderId} \neq 0 \,\}$\;
  $S_{\mathrm{font}} \leftarrow \{\, i \mid (xfs[i].\mathrm{fillId} \in \{0,\text{None}\}) \land (xfs[i].\mathrm{borderId} \in \{0,\text{None}\}) \land \neg xfs[i].\mathrm{applyFill} \land \neg xfs[i].\mathrm{applyBorder} \,\}$\;

  \tcp{3) Whitespace sharedStrings indices}
  $W \leftarrow \IdxWhitespaceSI(\text{sharedStringsXML})$\;

  \tcp{4) Bytes-keyed XF map (for s="...")}
  $\mathrm{xfsMap} \leftarrow \{\,\texttt{"0"}\mapsto xfs[0],\ \texttt{"1"}\mapsto xfs[1],\dots\}\;$

  \tcp{5) Consolidated branches that also build the regexes}

  %%% ONLY\_FONT\_STYLES + (rho1, rho2)
  \uIf{$|S_{\mathrm{font}}|>0$}{
    $\mathrm{ONLY\_FONT\_STYLES} \leftarrow \Join(S_{\mathrm{font}}, \texttt{|})$\;
    \tcp{irrelevant row attributes for empty rows pattern: $p_{\text{ra}}$}
    $p_{\text{ra}} \leftarrow$ \code{'\s*(\bspans="[^"]+?"|\b[^\s:]*?'} + \\
    \code{':dyDescent="[^"]*"|\bs="(?:'} \\
    \code{ + ONLY\_FONT\_STYLES + ')")'}\;
    \tcp{consecutive font-only styled empty cells pattern: $p_{\text{e}}$}
    $p_{\text{e}} \leftarrow$ \code{'(?:<c\b(?:(?![^>]*\bs=)[^>]*|[^>]*\bs="(?:'} \\
    \code{ + ONLY\_FONT\_STYLES + ')"[^>]*)/>)+'}\;
  }\Else{
    $p_{\text{ra}} \leftarrow$ \code{'\s*(\bspans="[^"]+?"|\b[^\s:]*?'} \\
    \code{ + ':dyDescent="[^"]*")'}\;
    $p_{\text{e}} \leftarrow$ \code{'(?:<c\b(?![^>]*\bs=)[^>]*/>)+'}\;
  }

  %%% FB + (rho3)
  \uIf{$|S_{\mathrm{fb}}|>0$}{
    $\mathrm{FB} \leftarrow \Join(S_{\mathrm{fb}}, \texttt{|})$\;
    \tcp{fill/border style or hidden attribute pattern: $p_{\text{fbh}}$}
    $p_{\text{fbh}} \leftarrow$ \code{'(?:s="(?:' + FB + ')"|"hidden="1")'}\;
  }\Else{
    $\mathrm{FB} \leftarrow \text{None}$\;
    $p_{\text{fbh}} \leftarrow$ \code{'"hidden="1"'}\;
  }

  %%% WS + (rho4, rho5)
  \uIf{$|W|>0$}{
      $\mathrm{WS} \leftarrow \Join(W, \texttt{|})$\;
    \tcp{whitespace value pattern: $p_{\text{ws}}$ (<v>k</v>) }
    $p_{\text{ws}} \leftarrow$ \code{'<v>(?:' + WS + ')</v>'}\;
    \tcp{whitespace cells pattern: $p_{\text{wc}}$ (<c t="s"><v>k</v></c>)}
    $p_{\text{wc}} \leftarrow$ \code{'(<c\b[^>]*?\bt="s"[^>]*?)><v>(?:'} \\
    \code{   + WS + ')</v></c>'}\;
  }\Else{
    $p_{\text{ws}} \leftarrow \text{None}$\;
    $p_{\text{wc}} \leftarrow \text{None}$\;
  }

  \tcp{6) Pattern compilation}
  $R \leftarrow \{\textbf{pra},\textbf{pe},\textbf{pfb},\textbf{pws},\textbf{pwc}\}$ with \textsf{compile} for all non-\texttt{None}\;


  \Return $\textsf{Patterns}(\mathrm{xfsMap}, R)$\;
}
\caption{Consolidated construction: compute alternations and build dependent regexes in the same branch.}
\label{alg:build-wb-patterns}
\end{algorithm}



These patterns parameterise all sheet and workbook transformations that follow, ensuring we preserve semantically relevant visual styling.



\newcommand{\xtag}[1]{\text{\scriptsize\texttt{#1}}} % needs amsmath (acmart loads it)
\begin{algorithm}
  \caption{Streaming extraction of sheet components}
  \label{alg:preprocess}
  \DontPrintSemicolon
  \SetKwFunction{ParseCols}{parseColDefs}
  \SetKwFunction{PreSD}{preProcessSheetData}
  \SetKwFunction{TrackMC}{trackMaxMerge}

  \SetKwFunction{Preprocess}{PreprocessSheetStreamwise}
  \SetKwProg{Fn}{Function}{:}{}

  \Fn{\Preprocess{sheetStream, patterns, chunkSize}}{
    $state \gets \texttt{IN\_WORKSHEET}$;\quad $buf \gets \epsilon$;\;
    \textbf{initialise accumulators: }
      $B_{\mathrm{preCols}},\ C_{\mathrm{cols}},\ B_{\mathrm{preSD}},\ S_{\mathrm{data}},\ B_{\mathrm{preMC}},\ M_{\mathrm{cells}},\ T_{\mathrm{tail}}$;\;
    $processedTags \gets \varnothing$;\quad $colDefs \gets \varnothing$;\quad $maxMerge \gets \bot$;\;

    $\mathcal{O} \gets \{\xtag{<cols>},\ \xtag{<sheetData>},\ \xtag{<mergeCells>},\ \xtag{</worksheet>}\}$;\;
    $\phi_{\mathrm{cols}}(x) \gets \textsf{parseColDefs}(x,\ patterns)$;\;
    $\phi_{\mathrm{sd}}(x) \gets \textsf{preProcessSheetData}(x,\ patterns)$;\;
    $\phi_{\mathrm{mc}}(x) \gets \textsf{trackMaxMerge}(x)$;\;

    \While{$state \neq \texttt{FOUND\_END}$}{
      $buf \gets \textsf{concat}\!\bigl(buf,\ \textsf{read}(sheetStream,\ chunkSize)\bigr)$;\;

      \uIf{$state = \texttt{IN\_WORKSHEET}$}{
        $remainingTags \gets \mathcal{O} \smallsetminus processedTags$\;
        $(tag,\, span) \gets \textsf{findNextTag}\bigl(buf,\ remainingTags\bigr)$;\;


        \uIf{\textsf{found}(tag)}{
          \lIf{$tag = \xtag{<cols>}$}{ $B_{\mathrm{preCols}} \gets \textsf{concat}\!\bigl(B_{\mathrm{preCols}},\ \textsf{takeUntil}(buf,\ tag)\bigr)$ }
          \lIf{$tag = \xtag{<sheetData>}$}{ $B_{\mathrm{preSD}} \gets \textsf{concat}\!\bigl(B_{\mathrm{preSD}},\ \textsf{takeUntil}(buf,\ tag)\bigr)$ }
          \lIf{$tag = \xtag{<mergeCells>}$}{ $B_{\mathrm{preMC}} \gets \textsf{concat}\!\bigl(B_{\mathrm{preMC}},\ \textsf{takeUntil}(buf,\ tag)\bigr)$ }

          $buf \gets \textsf{after}(buf,\ tag)$;\;
          $state \gets \textsf{nextState}(tag)$;\;
          $processedTags \gets processedTags \cup \{tag\}$;\;
        }
      }

      \uElseIf{$state = \texttt{IN\_COLS}$}{
        $(chunk,\ buf,\ state,\ colDefs) \gets$\\
        \Indp $\textsf{consumeUntil}\bigl(buf,\ \xtag{</cols>},\ \phi_{\mathrm{cols}}\bigr)$\Indm;\;
        \textsf{appendInPlace} \!($C_{\mathrm{cols}},\ chunk$)\;
      }

      \uElseIf{$state = \texttt{IN\_SHEETDATA}$}{
        $(chunk,\ buf,\ state) \gets$\\
        \Indp $\textsf{consumeUntil}\bigl(buf,\ \xtag{</sheetData>},\ \phi_{\mathrm{sd}}\bigr)$\Indm;\;
        \textsf{appendInPlace} \!($S_{\mathrm{data}},\ chunk$)\;
      }

      \uElseIf{$state = \texttt{IN\_MERGECELLS}$}{
        $(chunk,\ buf,\ state,\ maxMerge) \gets$\\
        \Indp $\textsf{consumeUntil}\bigl(buf,\ \xtag{</mergeCells>},\ \phi_{\mathrm{mc}}\bigr)$\Indm;\;
        \textsf{appendInPlace} \!($M_{\mathrm{cells}},\ chunk$)\;
      }

      \Else{
        \textsf{appendInPlace} \!($T_{\mathrm{tail}},\ \textsf{rest}(buf)$)\; \textbf{break}\;
      }
    }

\Return $\textsf{Preprocessed}\left(
  \vcenter{\hbox{$
    \begin{alignedat}{1}
      &B_{\mathrm{preCols}},\ C_{\mathrm{cols}},\ B_{\mathrm{preSD}},\ S_{\mathrm{data}},\\
      &B_{\mathrm{preMC}},\ M_{\mathrm{cells}},\ T_{\mathrm{tail}},\\
      &\textsf{Metrics}(colDefs,\ maxMerge)
    \end{alignedat}
  $}}
\right)$;\;

  }

  \BlankLine

  % --- helper functions (stacked) ---
  \Fn{\ParseCols{$x$, patterns}}{
    $defs \gets \textsf{newRangeMap}()$\;
    \ForEach{$c \gets \textsf{nextCol}(x)$}{
      $(\texttt{min},\texttt{max},\texttt{style},\texttt{hidden}) \gets \textsf{fields}(c)$\;
      $xf \gets \textsf{xfLookup}(\texttt{style},\ patterns)$\tcp*[r]{map raw index to style facts}
      \textsf{addRange}($defs$, $\texttt{min}$, $\texttt{max}$, $xf$, $\texttt{hidden}$)\;
    }
    \textsf{finalise}($defs$)\;
    \KwRet $defs$\;
  }


  \BlankLine

  \Fn{\TrackMC{$x$}}{
    $\mathcal{M} \gets \textsf{findAllMergeRefs}(x)$\tcp*[r]{each ref is $A\!:\!B$ with $(\mathrm{col}_1,\mathrm{row}_1)$–$(\mathrm{col}_2,\mathrm{row}_2)$}
    \If{$\mathcal{M}=\varnothing$}{\KwRet $\bot$}
    $\mathcal{C} \gets \varnothing$;\ $\mathcal{R} \gets \varnothing$\;
    \ForEach{$r \KwIn \mathcal{M}$}{
      $(c_1, r_1, c_2, r_2) \gets \textsf{splitRef}(r)$\;
      $\mathcal{C} \gets \mathcal{C} \cup \{c_1,c_2\}$;\quad $\mathcal{R} \gets \mathcal{R} \cup \{r_1,r_2\}$\;
    }
    $c^\star \gets \textsf{maxCol}(\mathcal{C})$;\quad $r^\star \gets \textsf{maxRow}(\mathcal{R})$\tcp*[r]{column order via A1 letters $\to$ index}
    \KwRet $\langle c^\star,\ r^\star\rangle$\;
  }

\end{algorithm}


\begin{algorithm}[t]
\caption{splitSheetDataRanges — compute \emph{cells\_range} and trailing \emph{rows\_range}}
\label{alg:split-sheetdata}
\DontPrintSemicolon
\SetKwProg{Fn}{Function}{:}{}

\Fn{splitSheetDataRanges($x$, $patterns$)}{
  \tcp{Input $x$: bytes of the \xtag{<sheetData>} node.}
  \tcp{$patterns.p_{\mathrm{fbhid}}$: compiled regex for \code{s="FB"} or \code{hidden="1"}.}

  \BlankLine
  \tcp{1) Find the end of the last row that still contains a cell.}
  $i_{\mathrm{lastC}} \gets \textsf{rfind}(x,\ \code{"<c "})$\;
  \If{$i_{\mathrm{lastC}} = -1$}{
    $i_{\mathrm{cellsEnd}} \gets -1$\;
  }
  \Else{
    $j \gets \textsf{find}(x,\ \code{"</row>"},\ i_{\mathrm{lastC}})$\;
    \If{$j=-1$}{\KwRet \textsf{error}(\code{"no </row> after last <c"})}
    $i_{\mathrm{cellsEnd}} \gets j + |\code{"</row>"}|$\;
  }

  \BlankLine
  \tcp{2) Scan for last visible attribute (fill/border style or \code{hidden="1"}).}
  $i_{\mathrm{visAttr}} \gets i_{\mathrm{cellsEnd}}$\;
  $m_{\mathrm{last}} \gets \bot$\;
  \ForEach{$m \in \textsf{finditer}(patterns.p_{\mathrm{fbhid}},\ x,\ \text{start}=i_{\mathrm{cellsEnd}})$}{
    $m_{\mathrm{last}} \gets m$\;
  }
  \If{$m_{\mathrm{last}} \neq \bot$}{
    $i_{\mathrm{visAttr}} \gets \max(i_{\mathrm{visAttr}},\ \textsf{start}(m_{\mathrm{last}}))$\;
  }

  \BlankLine
  \tcp{3) From that anchor, find the end position of the last visible row.}
  $j_{\mathrm{close}} \gets \textsf{find}(x,\ \code{"</row>"},\ i_{\mathrm{visAttr}})$\;
  $j_{\mathrm{self}}  \gets \textsf{find}(x,\ \code{"/>"},\ i_{\mathrm{visAttr}})$\;

  \uIf{$j_{\mathrm{close}} \neq -1 \land j_{\mathrm{self}} \neq -1$}{
    \uIf{$j_{\mathrm{close}} < j_{\mathrm{self}}$}{
      $i_{\mathrm{rowsEnd}} \gets j_{\mathrm{close}} + |\code{"</row>"}|$\;
    }\Else{
      $i_{\mathrm{rowsEnd}} \gets j_{\mathrm{self}} + |\code{"/>"}|$\;
    }
  }
  \uElseIf{$j_{\mathrm{close}} = -1 \land j_{\mathrm{self}} \neq -1$}{
    $i_{\mathrm{rowsEnd}} \gets j_{\mathrm{self}} + |\code{"/>"}|$\;
  }
  \uElseIf{$j_{\mathrm{close}} \neq -1 \land j_{\mathrm{self}} = -1$}{
    $i_{\mathrm{rowsEnd}} \gets j_{\mathrm{close}} + |\code{"</row>"}|$\;
  }
  \Else{
    \KwRet \textsf{error}(\code{"no </row> or /> after last visible anchor"})\;
  }

  \BlankLine
  \tcp{4) Slice out the two ranges.}
  \uIf{$i_{\mathrm{cellsEnd}} = -1$}{
    $\textit{cells\_range} \gets \epsilon$\;
  }\Else{
    $\textit{cells\_range} \gets x[0 : i_{\mathrm{cellsEnd}}]$\;
  }

  $i_{\mathrm{rowsStart}} \gets \max(0,\ i_{\mathrm{cellsEnd}})$\;
  \uIf{$i_{\mathrm{rowsEnd}} = -1$}{
    $\textit{rows\_range} \gets \epsilon$\;
  }\Else{
    $\textit{rows\_range} \gets x[i_{\mathrm{rowsStart}} : i_{\mathrm{rowsEnd}}]$\;
  }

  \BlankLine
  \KwRet $(\textit{cells\_range},\ \textit{rows\_range})$\;
}
\end{algorithm}

\subsection{Streaming Sheet Pre-processing}



We implement our spreadsheet parser streaming-based, processing the XML data incrementally to minimize memory consumption and permit early discarding of unnecessary file sections. Both our “consecutive” and “interleaved” parsing approaches, which reduce memory usage, use this same underlying state machine. Below, we focus on how the parser transitions between states (i.e., how each state function decides to proceed), deferring lower-level details of in-state processing until later sections.
The complete streaming pre-processing algorithm is summarised in Algorithm \ref{alg:preprocess}.

\paragraph{States Overview.}
We define six main states in our parser:
\begin{itemize}
  \item \texttt{BEFORE\_COLS}: The parser has not yet encountered the \texttt{<cols>} element. It searches for either \texttt{<cols>} or \texttt{<sheetData>}.
  \item \texttt{IN\_COLS}: The parser found \texttt{<cols>} and is processing column definitions.
  \item \texttt{BEFORE\_SHEET\_DATA}: The parser has finished reading (or did not detect) the \texttt{<cols>} node and is searching for the first \texttt{<sheetData>} element.
  \item \texttt{IN\_SHEET\_DATA}: The parser is processing the main worksheet data (rows and cells).
  \item \texttt{AFTER\_SHEET\_DATA}: The parser has detected the closing \texttt{</sheetData>} and is collecting any subsequent chunks until the end of the file.
  \item \texttt{END\_OF\_FILE}: The parser has reached the end of the stream, and no further parsing is needed.
\end{itemize}

\paragraph{In between nodes}
In \texttt{BEFORE\_COLS}, \texttt{BEFORE\_SHEET\_DATA}, or \texttt{AFTER\_SHEET\_DATA}, the parser looks for tag boundaries (e.g., \verb|<cols>|, \verb|<sheetData>|) or the end of file. When it detects such a tag at position \verb|i| in the buffer:
\begin{enumerate}
  \item It \emph{cuts} the buffer at \verb|i|, splitting the buffer into two segments:
   (1) from the start to the end of the located tag, and (2) the bytes after that tag.
  \item It returns the first segment as a chunk to add to the output, which the parser outputs or processes for the current state.
  \item The second segment becomes the updated buffer, ensuring the XML node boundary is never split across segments.
\end{enumerate}
If the parser fails to find any relevant tag or end of file, it returns no new output chunk, remains in the same state and leaves the buffer unmodified. This ensures that subsequent incoming data can be appended, eventually providing enough content to match the target boundary.

\paragraph{In column definition node.}
In \texttt{IN\_COLS}, the parser searches for \verb|</cols>| or any \verb|<col .../>| elements. On detecting \verb|</cols>|, it cuts the buffer right before the closing tag and transitions to \texttt{BEFORE\_SHEET\_DATA}. If it does not find the closing tag but locates complete \verb|<col .../>| elements, those are consumed and removed from the buffer as a clean cut. If no match is found, the buffer remains unchanged, waiting for more data to complete the matching.


\begin{algorithm}
\caption{preProcessSheetData — streaming prefilters for \texttt{<sheetData>} chunks}
\label{alg:presd}
\DontPrintSemicolon
\SetKwFunction{PreSD}{preProcessSheetData}
\SetKwFunction{Compile}{compile}
\SetKwProg{Fn}{Function}{:}{}

\Fn{\PreSD{$x$, $patterns$}}{
  \tcp{Input: $x$ = \texttt{<sheetData>} bytes (row-aligned).}
  \tcp{Params: $patterns$ = pattern pack with fields
  $\{p_{\text{ws}}, p_{\text{wc}}, p_{\text{e}}\}$: whitespace, whitespace-cell, consecutive-font-only-styled-empty-cells.}
  $y \gets x$\;

    
 

  \tcp{Local constant: collapse consecutive hidden rows (keep \textbackslash1).}
  {\small
  $\textit{hid\_src} \gets$ \;
  \Indp
    \code{'(<row\b[^/>]*?\bhidden="1"[^/>]++'} +
    \code{'(?:/>|>(?:(?!(?:</row>|<v>)).)++</row>))'} +
    \code{'(?:<row\b[^/>]*?\bhidden="1"[^/>]++'} +
    \code{'(?:/>|>(?:(?!(?:</row>|<v>)).)++</row>))+'} \;
  \Indm
  }%
  $p_{\text{hid}} \gets \Compile(\textit{hid\_src})$\;


  \BlankLine
  \tcp{Coalesce consecutive hidden rows.}
  \If{$\textsf{has}(y,\ \xtag{hidden="1"})$}{
    $y \gets \textsf{sub}(p_{\text{hid}}, "\textbackslash1", y)$\;
  }

  \BlankLine
  \tcp{Replace whitespace shared-string cells.}
  \If{$\textsf{has}(p_{\text{ws}})\ \land\ \textsf{has}(p_{\text{wc}})$}{
    \If{$\textsf{match}(p_{\text{ws}}, y)$}{
      $y \gets \textsf{sub}(p_{\text{wc}}, "\textbackslash1/>", y)$\;
    }
  }

  \BlankLine
  \tcp{Drop consecutive empty, layout-neutral cells.}
  \If{$\textsf{has}(y,\ \xtag{<c})$}{
    $y \gets \textsf{sub}(p_{\text{e}}, "", y)$\;
  }

  \BlankLine
  \tcp{Return cleaned chunk.}
  \KwRet $y$\;
}
\end{algorithm}



\paragraph{In sheet data node.}
While \texttt{IN\_SHEET\_DATA}, the parser similarly looks for \verb|</sheetData>| or row boundaries \texttt{<row ...>...</row>} (or \texttt{<row .../>}). If \verb|</sheetData>| is found, the parser cuts the buffer at that tag and transitions to \texttt{AFTER\_SHEET\_DATA}. If no boundary is present yet, the buffer remains unaltered to combine with future data chunks.

\paragraph{Ensuring ``Clean Cuts'' within the Streaming Buffer.}
Our parser reads raw XML data into a working buffer and ensures each XML node boundary is cleanly preserved, so no node is split across buffers. Each state-transition function adheres to this principle, only cutting when a relevant XML tag is fully identified. 

\paragraph{Buffer Updates and State Progress.}
By constraining parsing decisions to these well-defined cuts, the parser never splits an XML node across multiple buffers. Each function (1) cuts on the matching tag boundary and updates the parser state, or (2) leaves the buffer intact if no complete boundary is found, thus allowing the next chunk to complete it. This way the parser can process the XML incrementally without needing to store the entire document in memory and it is ensured that the regular expressions dont't fail due to incomplete XML nodes and also only the regular expressions that are relevant to the current section of the document are applied, incresing the performance in comparison to applying all regular expressions on the whole document.


\begin{algorithm}
\caption{cleanRowsAndCellsInRange — apply $p_{\text{row}}$ with
cell cleanup}
\label{alg:rows-cells}
\DontPrintSemicolon
\SetKwFunction{Sub}{sub}
\SetKwProg{Fn}{Function}{:}{}

\Fn{cleanRowsAndCellsInRange($x$, $patterns$, $metrics$, $params$)}{
  \tcp{Input $x$: the cells\_range bytes.}

  {\small
  $\textit{row_pattern} \gets$ \;
  \Indp
    \code{'<row\b((?:[^>]*\br="([^"]+?)")?'} +
    \code{'(?:[^>]*\bhidden="(1)")?'} +
    \code{'(?:[^>]*\bs="([^"]+?)")?'} +
    \code{'(?:[^>]*\bht="([^"]+?)")?'} +
    \code{'[^/>]*?)(?:/>|>(.*?)</row>)'} \;
  \Indm
  }%
  $patterns.p_{\text{row}} \gets \Compile(\textit{row\_pattern})$\;

{\small
  $\textit{cell_pattern} \gets$ \;
  \Indp
    \code{'<c\b((?:[^>]*\br="([A-Z]+)(\d+)")?'} +
    \code{'(?:[^>]*\bs="(\d+?)")?[^/>]*?)/>'} \;
  \Indm
  }%
  $patterns.p_{\text{cell}} \gets \Compile(\textit{cell_pattern})$\;      
  
  \tcp{$metrics$: mutable per-sheet state; $params$: xfs/columns.}
  \KwRet \Sub($x$, $p_{\text{row}}$,
    \textsf{ReplaceRow}$(\cdot, patterns, metrics, params)$, $params$)\;
}
\end{algorithm}


The sub-routine \textsf{preProcessSheetData} (Algorithm \ref{alg:presd}) performs three linear passes over chunks:
(i) collapse explicit whitespace cells discovered via shared-strings indices,
(ii) remove consecutive empty cells (either styleless or font-only),
and (iii) opportunistically remove runs of hidden rows
while retaining the first sentinel
when needed to preserve row height/visibility semantics.

Two of these transformations depend on workbook-specific patterns
derived from \texttt{styles.xml} and \texttt{sharedStrings.xml} as described in Section
\ref{sec:wb-patterns}. Those are the removal of explicit whitespace cells
and the removal of consecutive empty cells with font-only or no styling.
Before these, we can also collapse consecutive hidden rows
into a single sentinel row.
In Excel’s internal representation,
rows can be marked as hidden using the attribute \xtag{hidden="1"}.
While such rows are visually collapsed in the spreadsheet,
their presence in the XML can extend unnecessarily over thousands or even millions
of entries. In practice, we observed worksheets where, starting at a certain point,
all subsequent rows were marked as hidden.
However, Excel renders all following rows as hidden
as soon as a single row with \xtag{hidden="1"} is encountered,
unless a later row explicitly overrides this setting.
Therefore, from a rendering perspective,
a long sequence of hidden rows is semantically equivalent to just the first hidden row
being present.
To exploit this, we use a regular expression to identify
and collapse consecutive hidden rows into a single representative row.
This technique avoids expensive per-row checks and requires no document-wide context,
making it highly suitable for pre-processing.
The resulting simplification significantly reduces the size of the document
before it is fully loaded into memory.
The regex used for this task is constructed and applied in Algorithm \ref{alg:presd} first.
Then the whitespace cell replacement and empty cell removal follow.


%\subsection{5. Row/Cell Cleaning Semantics}
%
%
%
%Once all chunks have been pre-cleaned,
%they are reassembled into a complete document for the second phase.
%This phase applies context-aware operations,
%such as identifying the last visually relevant row or cell.
%If trailing rows beyond this point contain no visual content—e.g.,
%they have a defined height but no visible cells—they can be safely removed
%without affecting the worksheet layout.
%However, this step is deferred until the full document structure is available,
%as it requires global context.
%
%Ultimately, even though the full workbook must be loaded into memory for this final step
%(e.g., to determine the truncation point),
%the overall memory load is substantially reduced
%thanks to the prior removal of non-visible cells and redundant hidden rows.
%In contrast, more expensive operations such as regex-based pattern detection
%on row content are reserved for the post-processing phase
%and only applied when necessary.
%
%We define conservative rules to preserve visible meaning:
%\begin{itemize}
%\item A \emph{cell} without a style is removed; a cell with style that \emph{overrides} row/column fill/border is condensed to \emph{self-closing} (\texttt{<c .../>}); font-only styles are treated as non-layout and removable.
%\item A \emph{row} that is empty and neither hidden nor styled with fill/border is removed. Empty rows with height or with fill/border become self-closing (\texttt{<row .../>}) after stripping irrelevant attributes (\emph{spans}, vendor attributes). Consecutive hidden rows are coalesced to a single sentinel.
%\end{itemize}
%
%The cleaning procedure is summarised in Algorithm \ref{alg:rows}.
%It first checks if the row is empty
%and if it was already empty before cleaning.
%If the row is empty, it checks if it is hidden and if the previous row was also hidden.
%If both conditions are true, the row is deleted.
%If not, it replaces the row with a self-closing tag.
%If the row has no style or fill or border and no height, it deletes the row.
%If the row became empty by cleaning, it adds a custom height if absent.
%Finally, it replaces the row with a self-closing tag.



\begin{algorithm}
\caption{ReplaceRow — parse attributes, clean cells, keep/delete}
\label{alg:replace-row}
\DontPrintSemicolon
\SetKwFunction{CleanCells}{CleanCells}
\SetKwFunction{GenDelRow}{generateAndLogDeletedRow}
\SetKwFunction{GenSelfRow}{generateAndLogSelfClosingRow}
\SetKwFunction{GenRow}{generateAndLogRow}
\SetKwProg{Fn}{Function}{:}{}

\Fn{ReplaceRow($m$, $patterns$, $metrics$, $params$)}{
  \tcp{$m$ has groups from $p_{\text{row}}$:}
  \tcp{$rowAttr$, $r$, $hidden$?, $styleId$?, $ht$?, $content$?}
  $(rowAttr,\ r,\ hidden,\ styleId,\ ht,\ content)
    \gets \text{groups}(m)$\;
  $rowStyle \gets
    (styleId\neq\bot)\ ?\ params.\mathrm{xfsMap}[styleId]\ :\ \bot$\;
  $hasHeight \gets (ht\neq\bot)$,\quad $wasEmpty \gets (content=\epsilon)$\;

  \tcp{Persist row attributes for later checks.}
  $metrics.\text{row\_attributes}[r] \gets
    (rowAttr,\ r,\ (hidden=1),\ rowStyle,\ hasHeight,\ wasEmpty)$\;

  \tcp{Clean inner cells if present.}
  \If{$\neg wasEmpty$}{
    $content \gets \CleanCells(content,\ patterns,\ metrics,\ params)$\;
  }
  $isEmpty \gets (content=\epsilon)$,\quad
  $becameEmpty \gets (\neg wasEmpty \land isEmpty)$\;

  \tcp{Hidden rows: keep first of a run, drop subsequent.}
  \If{$hidden=1$}{
    \If{$\text{previousRowWasAlsoHidden}(metrics)$}{
      \KwRet \GenDelRow($metrics,\ (rowAttr,\ r,\ \cdot)$)\;
    }
    \Else{
      \KwRet \GenSelfRow($rowAttr,\ params,\ metrics,\ (rowAttr,\ r,\ \cdot)$)\;
    }
  }

  \tcp{Non-hidden: drop rows empty \& layout-neutral.}
  \If{$isEmpty$}{
    \If{$(rowStyle=\bot \lor \lnot rowStyle.\text{has\_fill\_or\_border})
         \land \lnot hasHeight$}{
      \KwRet \GenDelRow($metrics,\ (rowAttr,\ r,\ \cdot)$)\;
    }
    \If{$becameEmpty$}{
      $rowAttr \gets \textsf{addCustomHeightIfAbsent}(rowAttr)$\;
    }
    \KwRet \GenSelfRow($rowAttr,\ params,\ metrics,\ (rowAttr,\ r,\ \cdot)$)\;
  }

  \tcp{Keep non-empty rows with cleaned cells.}
  \KwRet \GenRow($rowAttr,\ content,\ metrics,\ (rowAttr,\ r,\ \cdot)$)\;
}
\end{algorithm}


\subsection{Post-processing of \texttt{<sheetData>}}
\label{sec:postproc}

\noindent
After streaming prefilters and structural range extraction, the post-processing
phase transforms the \texttt{<sheetData>} payload into a compact but
layout-equivalent representation. The phase comprises two linear passes that
operate on disjoint byte slices returned by the splitter of
Alg.~\ref{alg:split-sheetdata}:
(i) a \emph{rows+cells} sweep over the prefix slice that still contains cells
(\emph{cells\_range}), and (ii) a \emph{trailing empty-rows} sweep over the
suffix slice that no longer contains any cells (\emph{rows\_range}).
Both passes are guided by workbook-specific patterns (compiled as described in
\S\ref{sec:wb-patterns}) and update a small, per-sheet metrics state used by later
steps (merge-column normalization and diagnostics).



\code{metrics} (kept/deleted rows/cells, row attributes, etc.).


\paragraph*{Range splitting.}

Alg.~\ref{alg:split-sheetdata} identifies the last row that still contains a
cell by scanning backwards for a \xtag{<c} token, then advancing to the
end of that row. From that anchor it scans forward for the last
\emph{visibly relevant attribute}---either a fill/border-bearing style
(\code{s="FB"}) or \xtag{hidden="1"}---by iterating a precompiled pattern
\mbox{$p_{\mathrm{fbhid}}$}. The earliest of the next
\xtag{</row>} and \xtag{/>} delimiters after this anchor marks the end of the
trailing visible-rows region. The function returns two slices: the \emph{cells_range} from the beginning up to and including
the last cell-bearing row, and the \emph{rows_range} from there up to the last
visibly relevant row. This design avoids a full parse and guarantees
$O(|x|)$ time with a constant number of linear scans over bytes.

\paragraph*{Row+cell sweep on the cells prefix.}
The first pass (Alg.~\ref{alg:rows-cells}) applies a row-level regex
\mbox{$p_{\text{row}}$} to the \emph{cells\_range}. Each row match is handed to
\emph{ReplaceRow} (Alg.~\ref{alg:replace-row}), which (i) records row
attributes for later hidden/override checks, (ii) cleans the inner
cells via \emph{CleanCells} (Alg.~\ref{alg:clean-cells}) if present,
and (iii) decides to delete the row, keep it as self-closing (empty but
layout-relevant), or keep it with (cleaned) content. The cell
policy is encapsulated in \emph{ReplaceCell} (Alg.~\ref{alg:replace-cell}):
cells in hidden row/column, without style, or without visible effects are
dropped; cells with fills/borders or with effective overrides are kept as
self-closing. All actions are appended to \code{metrics} for downstream use.



\paragraph*{Trailing empty-rows sweep.}
The second pass (Alg.~\ref{alg:rows-trailing}) runs over the \emph{rows\_range}
using \mbox{$p_{\text{rowE}}$}. \emph{ReplaceEmptyRow}
(Alg.~\ref{alg:replace-empty-row}) preserves the first row of any hidden-run and
any style-bearing empty row (to retain layout height/borders), while deleting
layout-neutral empty rows. This pass never touches cells (by construction there
are none in the suffix), which makes it both simpler and faster.

\begin{algorithm}[h]
\caption{cleanEmptyRowsFromAfterCellRange — apply $p_{\text{rowE}}$}
\label{alg:rows-trailing}
\DontPrintSemicolon
\SetKwFunction{Sub}{sub}
\SetKwProg{Fn}{Function}{:}{}

\Fn{cleanEmptyRowsFromAfterCellRange($x$, $patterns$, $metrics$, $params$)}{
  \tcp{Input $x$: rows\_range bytes (after last row with any cell).}
  \KwRet \Sub($x$, $p_{\text{rowE}}$,
    \textsf{ReplaceEmptyRow}$(\cdot, metrics, params)$, $params$)\;
}
\end{algorithm}


\begin{algorithm}[h]
\caption{ReplaceEmptyRow — keep only hidden-first or style rows}
\label{alg:replace-empty-row}
\DontPrintSemicolon
\SetKwFunction{GenDelRow}{generateAndLogDeletedRow}
\SetKwFunction{GenSelfRow}{generateAndLogSelfClosingRow}
\SetKwProg{Fn}{Function}{:}{}

\Fn{ReplaceEmptyRow($m$, $metrics$, $params$)}{
  \tcp{Groups from $p_{\text{rowE}}$: $rowAttr$, $r$, $hidden$?, $styleId$?}
  $(rowAttr,\ r,\ hidden,\ styleId) \gets \text{groups}(m)$\;
  $rowStyle \gets
    (styleId\neq\bot)\ ?\ params.\mathrm{xfsMap}[styleId]\ :\ \bot$\;

  $metrics.\text{row\_attributes}[r] \gets
    (rowAttr,\ r,\ (hidden=1),\ rowStyle,\ \text{hasHeight}=false,\ \text{isEmpty}=true)$\;

  \If{$hidden=1$}{
    \If{$\text{previousRowWasAlsoHidden}(metrics)$}{
      \KwRet \GenDelRow($metrics,\ (\cdot)$)\;
    }
    \Else{
      \KwRet \GenSelfRow($rowAttr,\ params,\ metrics,\ (\cdot)$)\;
    }
  }

  \If{$(rowStyle=\bot \lor \lnot rowStyle.\text{has\_fill\_or\_border})$}{
    \KwRet \GenDelRow($metrics,\ (\cdot)$)\;
  }
  \Else{
    \KwRet \GenSelfRow($rowAttr,\ params,\ metrics,\ (\cdot)$)\;
  }
}
\end{algorithm}

 
\paragraph*{Metrics and column merging.}
The rows+cells pass records kept cell coordinates and row attributes in
\code{metrics}; combined with merge-cell bounds and visible column definitions,
\code{get_max_column} computes the global rightmost used column. This index is
then used to merge extraneous \xtag{<col>} ranges, as outlined in
\S\ref{sec:columns}.

 
















% ---------------------------------------------------------------
% Helper: cell replacement policy (decision summary)
% ---------------------------------------------------------------
\paragraph{Cell policy (summary).}
A cell is kept (as self-closing) iff it visibly affects layout:
not in a hidden row/column; has a style; either carries a
fill/border or overrides an \emph{apply} style from its row/column.
Otherwise the cell is deleted.


% ---------------------------------------------------------------
% CleanCells: apply p_cell using ReplaceCell policy
% ---------------------------------------------------------------
\begin{algorithm}[t]
\caption{CleanCells — apply $p_{\text{cell}}$ with ReplaceCell}
\label{alg:clean-cells}
\DontPrintSemicolon
\SetKwFunction{Sub}{sub}
\SetKwProg{Fn}{Function}{:}{}

\Fn{CleanCells($rowBytes$, $patterns$, $metrics$, $params$)}{
  \If{$rowBytes=\epsilon$}{\KwRet $rowBytes$}
  \KwRet \Sub($rowBytes$, $p_{\text{cell}}$,
    \textsf{ReplaceCell}$(\cdot, metrics, params)$, $params$)\;
}
\end{algorithm}

% ---------------------------------------------------------------
% ReplaceCell: keep (self-close) or delete
% ---------------------------------------------------------------
\begin{algorithm}[t]
\caption{ReplaceCell — drop or keep self-closing cell}
\label{alg:replace-cell}
\DontPrintSemicolon
\SetKwFunction{GenDelCell}{generateAndLogDeletedCell}
\SetKwFunction{GenSelfCell}{generateAndLogSelfClosingCell}
\SetKwProg{Fn}{Function}{:}{}

\Fn{ReplaceCell($m$, $metrics$, $params$)}{
  \tcp{Groups from $p_{\text{cell}}$: $cellAttr$, $col$?, $row$?, $styleId$?}
  $(cellAttr,\ col,\ row,\ styleId) \gets \text{groups}(m)$\;
  $style \gets (styleId\neq\bot)\ ?\ params.\mathrm{xfsMap}[styleId]\ :\ \bot$\;
  $coord \gets (col,row)$\;

  \If{$\textsf{cellIsHidden}(coord,\ metrics)$}{\KwRet \GenDelCell($metrics,\ (\cdot)$)}
  \If{$style=\bot$}{\KwRet \GenDelCell($metrics,\ (\cdot)$)}
  \If{$style.\text{has\_fill\_or\_border}$}{\KwRet \GenSelfCell($metrics,\ (\cdot)$)}
  \If{$\lnot style.\text{has\_apply\_fill\_or\_border}$}{\KwRet \GenDelCell($metrics,\ (\cdot)$)}

  \If{$\textsf{cellOverridesFillOrBorder}(style,\ coord,\ metrics)$}{
    \KwRet \GenSelfCell($metrics,\ (\cdot)$)
  }
  \Else{
    \KwRet \GenDelCell($metrics,\ (\cdot)$)
  }
}
\end{algorithm}



% ---------------------------------------------------------------
% Notes: metrics interactions and column merging
% ---------------------------------------------------------------
\paragraph{Notes.}
The rows+cells pass accumulates kept cell coordinates in
\code{metrics.kept_cell_states}; combined with merge-cells bounds
and visible column definitions, \code{get_max_column} yields the
overall maximum column used to merge extraneous \xtag{<col>} ranges.
Both passes update \code{kept_row_states}, \code{deleted_row_states},
and the \code{row_attributes} map, which drive checks later in the pipeline.
% ---------------------- END DROP-IN ----------------------


%The cell rule references row/column styles through \emph{metrics} to detect overrides of fill/border;
%overridden cells are retained in self-closing form, otherwise removed.

\subsection{Column Range Compaction}
\label{sec:columns}

We compact \xtag{<cols>} by merging “extraneous” trailing definitions
that exceed the rightmost visible column inferred from (i) kept cells,
(ii) \xtag{<mergeCells>}, and (iii) last visible column definition (hidden or with fill/border).
If the first extraneous range precedes others, we widen it to the last, replacing the suffix.

\begin{algorithm}
\DontPrintSemicolon
\SetKwFunction{MergeCols}{MergeExtraneousColumns}
\SetKwProg{Fn}{Function}{:}{}
\Fn{\MergeCols{$C_{\text{cols}}$, maxCol}}{
\lIf{$maxCol=\bot$}{\Return $C_{\text{cols}}$}
$(p,,head) \leftarrow \textsf{scanUntilFirst}(\text{range.max} > maxCol)$;
\lIf{$p=\bot$}{\Return $C_{\text{cols}}$}
$last \leftarrow \textsf{scanLastRange}()$;
\lIf{$last=\bot$}{\Return $C_{\text{cols}}$}
$p.\text{max} \leftarrow last.\text{max}$;
\Return $head \mathbin{|} \textsf{serialise}(p)$
}
\caption{Compacting column ranges beyond the last visible column.}
\label{alg:cols}
\end{algorithm}

\subsection{7. Workbook De-duplication of Defined Names}

Some workbooks redundantly define both \texttt{Print\_Area} and \texttt{\_xlnm.Print\_Area} (and analogously \texttt{Print\_Titles}) for the same \texttt{localSheetId}. We keep one representative and delete duplicates.

\begin{algorithm}
\DontPrintSemicolon
\SetKwFunction{CleanWB}{CleanWorkbookDefinedNames}
\SetKwProg{Fn}{Function}{:}{}
\Fn{\CleanWB{workbookXML}}{
  \(D \gets \textsf{extractBlock}(workbookXML,\ \texttt{<definedNames>})\);\;
  \uIf{\(D=\bot\)}{\Return workbookXML}
  \tcp*[l]{Group duplicates by $(\textsf{name},\ \textsf{localSheetId})$}
  \ForEach{sheet \(s\)}{
    \If{\texttt{Print\_Area} \textbf{and} \texttt{\_xlnm.Print\_Area} exist for \(s\)}{delete one\;}
    \If{\texttt{Print\_Titles} \textbf{and} \texttt{\_xlnm.Print\_Titles} exist for \(s\)}{delete one\;}
  }
  \Return \textsf{replaceBlock}(workbookXML,\ D)\;
}
\caption{Workbook de-duplication of \texttt{<definedNames>}.}
\label{alg:wb}
\end{algorithm}



\subsection{10. Multi-process Execution}

We support multi-process execution with a process pool. We first analyse files to create work units (workbook task + per-sheet tasks) and then schedule \emph{clean} jobs. A per-file writer is lazily initialised and closed once all scheduled tasks for that file complete.

\begin{algorithm}
\DontPrintSemicolon
\SetKwFunction{RunMP}{RunMultiProcess}
\SetKwProg{Fn}{Function}{:}{}
\Fn{\RunMP{files, outPath, chunkSize, sheetFilter, onlyWB, \(W\)}}{
  \(pool \gets \textsf{procPool}(W)\);\ \(state \gets \varnothing\);\tcp*[l]{future \(\mapsto\) (task, file)}
  \ForEach{\(f \in files\)}{
    \(zipWriter[f] \gets \bot\);\ \((pat,\ wbTask,\ sheetTasks,\ others) \gets \textsf{plan}(f,\ sheetFilter,\ onlyWB)\);\;
    \(F \gets \textsf{submit}(pool,\ \textsf{CleanWorkbookDefinedNames},\ wbTask)\);\ \(state[F] \gets (wbTask,\ f,\ pat)\);\;
    \ForEach{\(s \in sheetTasks\)}{
      \(F \gets \textsf{submit}(pool,\ \textsf{CleanSheet},\ s,\ pat,\ chunkSize)\);\ \(state[F] \gets (s,\ f,\ pat)\);\;
    }
    \(pendingOthers[f] \gets others\);\;
  }
  \While{\(state \neq \varnothing\)}{
    \(F \gets \textsf{takeCompleted}()\);\ \((task,\ f,\ pat) \gets state[F]\);\ \textsf{del}\ \(state[F]\);\;
    \If{\(zipWriter[f]=\bot\)}{\(zipWriter[f]\gets \textsf{openZipForWrite}(f,\ outPath)\)}
    \textsf{write}\((zipWriter[f],\ task.name,\ \textsf{result}(F))\);\;
    \If{\textsf{allScheduledFor}\((f)\) \textbf{are written}}{
      \ForEach{\(o \in pendingOthers[f]\)}{\textsf{copy}\((zipWriter[f],\ o)\);\;}
      \textsf{close}\((zipWriter[f])\);\ \textsf{del}\ \(zipWriter[f]\);\;
    }
  }
}
\caption{Multi-process orchestration with lazy writers.}
\label{alg:multi}
\end{algorithm}

 






